\subsection{Ordinary eighth moments for every short-one terminal}\label{proof:Q5-C058-ordinary}
Let \(K\) be a real antisymmetric kernel with \(|K|\le\allowbreak 1\) on a nonempty finite probability space \((X,\allowbreak \mu)\). Use the finite notation, square envelope and signed coefficients of Section~\ref{proof:Q5-C058}, in particular \eqref{eq:L58:finiteinputs}, \eqref{eq:L58:propagation} and \eqref{eq:L58:receiving}. In this calculation write \(B_0=\allowbreak -K\) and \(g'=\allowbreak B_0\one=\allowbreak -g\). Put
\[
F_b=qf_b,\qquad X_{b,L}=T^L(qf_b),\qquad H_{b,L}=pX_{b,L}.
\tag{Q58O.0.1}\label{eq:Q58O:0.1}
\]
The terminal unit arm points inward, the \(b\)-edge arm points outward, every connector edge points toward the receiving root, and the receiving unit arm points outward. These directions are fixed by the original structure.

We prove the ordinary-eighth bound for every \(b\ge\allowbreak 3,\allowbreak L\ge\allowbreak 2\), with deficits \(d/2000\) for \(L\ge\allowbreak 4\), \(d/12000\) for \(L=\allowbreak 3\), and \(d/3000\) for \(L=\allowbreak 2\).
The \(b=\allowbreak 2,\allowbreak L\ge\allowbreak 2\) rows follow from Section~\ref{proof:Q5-C050}; the \(L=\allowbreak 1\) rows for every \(b\ge\allowbreak 2\) follow from Section~\ref{proof:Q5-C051}, using the scalar estimate of Section~\ref{proof:Q5-C049}. Their combination gives
\[
\boxed{\E H_{b,L}^8\le1-d/12000\qquad(b\ge2,\ L\ge1).}
\tag{Q58O.0.2}\label{eq:Q58O:0.2}
\]
All calculations below use \(\mu\). At \(d=\allowbreak 0\), leaf integration gives every message equal to one; divisions by \(d\) are used only for \(d>\allowbreak 0\).

\paragraph{Q58O.2 Fixed-node Hermite comparison and the Lucas polynomial}
Suppose \(N\) is real with \(|N|\le\allowbreak  u\), \(\E N\le\allowbreak -a\) and
\(\E N^2\le\allowbreak  hu^2\), where \(0<\allowbreak h\le\allowbreak 1/2\) and \(a\ge\allowbreak 0\).
Take \(c>\allowbreak 0\) and \(x=\allowbreak u/(2c)<\allowbreak 1\).
For \(f(z)=\allowbreak (1+z/(2c))^{16}\), quadratic Hermite interpolation with double contact at \(-hu\) and single contact at \(u\) is an upper bound on \([-u,\allowbreak u]\), since \(f'''\ge\allowbreak 0\) and \(z-u\le\allowbreak 0\).
Both the quadratic and the mean coefficients are nonnegative. Consequently
\[
E(1+N/(2c))^{16}\le M_h(x)-\frac a uD_h(x),
\tag{Q58O.2.1}\label{eq:Q58O:2.1}
\]

Here
\[
\begin{aligned}
M_h(x)&=\frac{h(1+x)^{16}+(1-hx)^{16}}{1+h},\\
D_h(x)&=\frac{1-h}{1+h}\,16x(1-hx)^{15}
+\frac{2h}{(1+h)^2}\bigl[(1+x)^{16}-(1-hx)^{16}\bigr].
\end{aligned}
\]
In this scalar comparison, write the degree-\(j\) coefficients as
\(\binom{16}{j}M_j\) and \(\binom{16}{j}D_j\). Then
\[
M_j=\frac{h+(-h)^j}{1+h},\quad
D_j=\frac{1-h}{1+h}j(-h)^{j-1}
+\frac{2h}{(1+h)^2}(1-(-h)^j).
\tag{Q58O.2.2}\label{eq:Q58O:2.2}
\]

We have \(M_0=\allowbreak 1\), \(M_1=\allowbreak 0\), \(M_2=\allowbreak h\), \(M_3=\allowbreak h(1-h)\), and
\(M_j\le\allowbreak  M_4=\allowbreak h(1-h+h^2)\) for \(j\ge\allowbreak 4\).
Every \(D_j\ge\allowbreak 0\). For even \(j\ge\allowbreak 4\), after factoring
\(h(1-h)/(1+h)^2\), the remainder is
\[
2\sum_{i=0}^{j-1}h^i-jh^{j-2}(1+h)>0,
\]
because \(jh^{j-2}\le\allowbreak 1\).
The low-degree values used below are
\[
D_1=1,\ D_2=0,\ D_3=h(2-h),\ D_4=2h(1-h)^2,
\quad D_5=2h-4h^2+6h^3-3h^4.
\tag{Q58O.2.3}\label{eq:Q58O:2.3}
\]


If \(c=\allowbreak 1-Ad\), \(u^2=\allowbreak Md\) and \(w=\allowbreak 1/c\), then
\(w^2-w=\allowbreak k_0x^2\), where \(k_0=\allowbreak 4A/M\).
Put \(v=\allowbreak w-1\ge\allowbreak 0\) and \(t=\allowbreak k_0x^2=\allowbreak wv\). The finite Lucas identity gives
\[
w^{16}+v^{16}=1+16t+104t^2+352t^3+660t^4
+672t^5+336t^6+64t^7+2t^8.
\]

Since \(v^{16}\le\allowbreak  t^8\),
\[
w^{16}\ge1+16t+104t^2+352t^3+660t^4
+672t^5+336t^6+64t^7+t^8.
\tag{Q58O.2.4}\label{eq:Q58O:2.4}
\]

The identity follows from \(S_0=\allowbreak 2\), \(S_1=\allowbreak 1\) and
\(S_n=\allowbreak S_{n-1}+tS_{n-2}\).

If \(\E N\le\allowbreak -b(d)d\), then
\(a/u\ge\allowbreak [2b(d)c/M]x\). Denote the coefficient in brackets by \(\ell(d)\).
Equations \eqref{eq:Q58O:2.1} and \eqref{eq:Q58O:2.4} reduce the estimate to a positive lower bound for
\(w^{16}-M_h(x)+\ell(d)xD_h(x)\), using the mean, second moment and row bound of the same \(N\).

\paragraph{Q58O.3 Three scalar estimates on the full continuous domains}
\emph{Scalar I.}
Assume \(|N|^2\le\allowbreak  d\), \(\E N\le\allowbreak -d/2\), \(\E N^2\le\allowbreak \theta d\), and set \(c=\allowbreak 1-3d/4\).
Then \(h=\allowbreak \theta\), \(k_0=\allowbreak 3\), \(c\ge\allowbreak 3/4\), \(x<\allowbreak 1\), and
\(\ell(d)=\allowbreak c\ge\allowbreak 3/4\).
Also \(M_4<\allowbreak 5/16\), \(560M_3<\allowbreak 136\), and \(560D_3>\allowbreak 350\).
After division by \(x^2\), the difference from \eqref{eq:Q58O:2.4} is at least
\[
\begin{aligned}
Q_I(x)&=45/4-136x+629x^2-1365x^3+7000x^4\\
&\quad -3575x^5+49000x^6-3575x^7+160000x^8\\
&\quad -1365x^9+244000x^{10}-175x^{11}+139000x^{12}-5x^{13}+6500x^{14}.
\end{aligned}
\]

The quartic coefficient is at least
\(936-1820(5/16)+(3/4)350=\allowbreak 629.75\).
From degree six onward, retain the positive Lucas terms and use \(M_j\le\allowbreak 5/16\).
The polynomial decomposes as
\[
\begin{aligned}
Q_I={}&(45/4-136x+420x^2)+x^2(209-1365x+2400x^2)\\
&+x^4(4600-3575x+1000x^2)+x^6(48000-3575x+100x^2)\\
&+x^8(159900-1365x+10x^2)+x^{10}(243990-175x+x^2)\\
&+x^{12}(138999-5x+x^2)+6499x^{14}.
\end{aligned}
\tag{Q58O.3.1}\label{eq:Q58O:3.1}
\]

Every bracket has negative discriminant; the first has minimum \(101/420\).
Thus every nonnegative \(H\le\allowbreak (c+N/2)^2\) satisfies
\[
EH^8\le1-(101/1680)c^{14}d\le1-d/2000.
\tag{Q58O.3.2}\label{eq:Q58O:3.2}
\]

Here \((3/4)^{14}>\allowbreak 1/64\).

\emph{Scalar II.}
Let \(M=\allowbreak 71/63\), and suppose
\[
|N|^2\le Md,\qquad \E N^2\le\theta Md,\qquad
\E N\le-(4/5-6d/5)d,\qquad c=1-4d/5.
\]
Then \(k_0=\allowbreak 1008/355>\allowbreak 17/6\), \(c\ge\allowbreak 11/15\), \(x<\allowbreak 1\), and
\[
\ell(d)=\frac{2c(4/5-6d/5)}M\ge\frac{44}{75M}>13/25.
\]

Since \(d=\allowbreak 4c^2x^2/M\le\allowbreak 4x^2/M\), we also have
\[
\ell(d)\ge\frac8{5M}-\frac{368}{25M^2}x^2.
\tag{Q58O.3.3}\label{eq:Q58O:3.3}
\]

Use \eqref{eq:Q58O:3.3} for \(D_1\), and \(\ell>\allowbreak 13/25\) for
\(D_3,\allowbreak D_4,\allowbreak D_5\); discard the remaining positive terms.
The difference divided by \(x^2\) is at least
\[
\begin{aligned}
Q_{II}(x)&=155/8-(1351/10)x+274x^2-960x^3+6500x^4\\
&\quad -3575x^5+38000x^6-3575x^7+120000x^8\\
&\quad -1365x^9+170000x^{10}-175x^{11}+90000x^{12}-5x^{13}+4000x^{14}.
\end{aligned}
\]

The quadratic coefficient is
\(384/(5M)-195/4>\allowbreak 155/8\), and
\(560M_3=\allowbreak 8645/64<\allowbreak 1351/10\).
The quartic coefficient is at least
\[
20736/(25M^2)-1820M_4+560D_3\cdot44/(75M)>274,
\]
using \(M^2<\allowbreak 1271/1000\), \(20736/(25M^2)>\allowbreak 652\),
\(M_4<\allowbreak 31/100\), and \(D_3>\allowbreak 16/25\).
The magnitude of the negative quintic coefficient is below
\[
4368(9/32)-1820(2/7)(13/25)=958.1<960,
\]
because \(M_5<\allowbreak 9/32\) and \(D_4>\allowbreak 2/7\).
The sextic coefficient exceeds \(8000-2483+1067=\allowbreak 6584>\allowbreak 6500\), by
\(352k_0^3>\allowbreak 8000\), \(M_6<\allowbreak 31/100\), and
\(D_5=\allowbreak 495885/1048576>\allowbreak 47/100\).
The remaining even coefficients, using \(k_0>\allowbreak 17/6\), exceed
\(38000,\allowbreak 120000,\allowbreak 170000,\allowbreak 90000,\allowbreak 4000\), respectively.
The magnitudes of the remaining negative odd coefficients, using
\(M_j<\allowbreak 31/100\), are below \(3575,\allowbreak 3575,\allowbreak 1365,\allowbreak 175,\allowbreak 5\).

The decomposition is
\[
\begin{aligned}
Q_{II}={}&(155/8-(1351/10)x+237x^2)+x^2(37-960x+6400x^2)\\
&+x^4(100-3575x+32000x^2)+x^6(6000-3575x+600x^2)\\
&+x^8(119400-1365x+10x^2)+x^{10}(169990-175x+x^2)\\
&+x^{12}(89999-5x+x^2)+3999x^{14}.
\end{aligned}
\tag{Q58O.3.4}\label{eq:Q58O:3.4}
\]

Every bracket is positive. The first has minimum
\(11549/94800>\allowbreak 1/10\), and the second has minimum one.
For \(0\le\allowbreak  H\le\allowbreak (c+N/2)^2\), this gives
\[
EH^8\le1-(M/40)c^{14}d\le1-d/12000.
\tag{Q58O.3.5}\label{eq:Q58O:3.5}
\]

Use \(c\ge\allowbreak 11/15>\allowbreak 2/3\) and \((2/3)^{14}>\allowbreak 1/300\).

\emph{Scalar III.}
Assume
\[
|N|^2\le d,\qquad \E N\le-41d/96,\qquad
\E N^2\le3d/10,\qquad c=1-37d/64.
\]
Then \(h=\allowbreak 3/10\), \(k_0=\allowbreak 37/16\), \(c\ge\allowbreak 155/192\), \(x<\allowbreak 1\), and
\(\ell(d)\ge\allowbreak 6355/9216>\allowbreak 11/16\).
The difference divided by \(x^2\) is at least
\[
\begin{aligned}
Q_{III}(x)&=12-118x+321x^2-632x^3+3500x^4\\
&\quad -2720x^5+15800x^6-2720x^7+42000x^8\\
&\quad -1040x^9+50000x^{10}-133x^{11}+22000x^{12}-4x^{13}+800x^{14}.
\end{aligned}
\]

Here
\[
\begin{gathered}
M_3=21/100,\quad M_4=237/1000,\quad M_5=2289/10000,\quad\\
D_3=51/100,\quad D_4=147/500,\quad D_5=3777/10000>3/8.
\end{gathered}
\]
The quadratic coefficient exceeds \(37-36+11=\allowbreak 12\).
The quartic coefficient exceeds
\(17797/32-431.34+196.35=\allowbreak 321.16625\).
The magnitude of the negative quintic coefficient is below
\[
4368(2289/10000)-(11/16)1820(147/500)=631.9677<632.
\]
The sextic coefficient exceeds \(557183/128-1898+1126>\allowbreak 3500\).
Using \(k_0^4>\allowbreak 200/7\) and \(k_0^2>\allowbreak 53/10\), the subsequent even coefficients exceed
\(15800,\allowbreak 42000,\allowbreak 50000,\allowbreak 22000,\allowbreak 800\); the remaining odd terms use
\(M_j\le\allowbreak 237/1000\).

We have the decomposition
\[
\begin{aligned}
Q_{III}={}&(12-118x+291x^2)+x^2(30-632x+3330x^2)\\
&+x^4(170-2720x+11000x^2)+x^6(4800-2720x+400x^2)\\
&+x^8(41600-1040x+10x^2)+x^{10}(49990-133x+x^2)\\
&+x^{12}(21999-4x+x^2)+799x^{14}.
\end{aligned}
\tag{Q58O.3.6}\label{eq:Q58O:3.6}
\]

Every bracket is positive, and the first has minimum \(11/291\).
Thus \(0\le\allowbreak  H\le\allowbreak (c+N/2)^2\) implies
\[
EH^8\le1-(11/1164)c^{14}d\le1-d/3000.
\tag{Q58O.3.7}\label{eq:Q58O:3.7}
\]

Use \(c>\allowbreak 4/5\) and \((4/5)^{14}>\allowbreak 1/25\).
The displayed terminating decimals denote exact finite rational numbers.

\paragraph{Q58O.4 Every $b\ge\allowbreak 3$ and $L\ge\allowbreak 4$}
In the finite envelope \eqref{eq:L58:propagation}, retain
\(\varepsilon_L\bar f_0\) and \(\varepsilon_{L-1}\bar f_1\).
Both coefficients are at least \(d/2\).
Moreover, \(R_L=\allowbreak V_L-q\) satisfies
\[
\E R_L\le-\tau\le-d/2,\qquad |R_L|^2\le d,\qquad
\E R_L^2\le\theta d.
\]
Hence
\[
X_{b,L}\le m^2[(1-d/2)q+R_L-d/2].
\]

From \(\beta_b\le\allowbreak  d\) and \(\beta_b\ge\allowbreak \kappa d\),
\[
m^2=1-\beta_b+\beta_b^2/4\le1-(11\kappa/12)d\le1-d/2.
\]
AM--GM gives
\[
H_{b,L}\le(1-3d/4+R_L/2)^2.
\tag{Q58O.4.1}\label{eq:Q58O:4.1}
\]

The square base is the average of two nonnegative factors, and
\(m^2/(1-d/2)\le\allowbreak 1\). Scalar I now gives
\(\E H_{b,\allowbreak L}^8\le\allowbreak 1-d/2000\).

\paragraph{Q58O.5 Every $b\ge\allowbreak 3$ at $L=\allowbreak 3$}
Expanding \(\bar f_2=\allowbreak q+r\) in \eqref{eq:L58:propagation} gives
\[
X_{b,3}\le m^2[V_3-\gamma r-a-B_*q],\quad a=d-k\rho\ge41d/80.
\]

Put \(N=\allowbreak V_3-\gamma r-q\). Its row polynomial is \(P_2-\gamma\).
The real part of \(P_2\) is at least \(1-(6/5)\theta^2>\allowbreak 4/5\), while
\(\gamma\le\allowbreak 2/5\); subtracting \(\gamma\) therefore does not increase the modulus.
Thus \(|N|^2\le\allowbreak  Md\) and \(\E N^2\le\allowbreak \theta Md\), where \(M=\allowbreak 71/63\).
Also
\[
-EN=d-kD_3-\gamma d
\ge(4/5-6d/5)d.
\tag{Q58O.5.1}\label{eq:Q58O:5.1}
\]

Here \(D_3\le\allowbreak \delta_3\le\allowbreak \theta^2d\) and
\(1-(6/5)\theta^2>\allowbreak 4/5\).
Let \(B=\allowbreak 1-(209/384)d\). The bound above gives \(m^2\le\allowbreak  B\), and
\(B_*\ge\allowbreak 5d/8>\allowbreak (209/384)d\).
The positive envelope and AM--GM give
\[
H_{b,3}\le[B-(41/160)d+N/2]^2
\le(1-4d/5+N/2)^2,
\tag{Q58O.5.2}\label{eq:Q58O:5.2}
\]

Indeed \(209/384+41/160=\allowbreak 1537/1920>\allowbreak 4/5\).
Scalar II proves \(\E H_{b,\allowbreak 3}^8\le\allowbreak 1-d/12000\).

\paragraph{Q58O.6 Every $b\ge\allowbreak 3$ at $L=\allowbreak 2$}
Keep the dependence between \(m,\allowbreak \gamma\) and the finite arm tail.
Let \(B=\allowbreak 1-\gamma>\allowbreak 0\) be a scalar, distinct from the operator \(B_0=\allowbreak -K\), and put \(R_2=\allowbreak r+kK^2R\).
The envelope is
\[
X_{b,2}\le m^2[Bq+R_2-a_1].
\]

AM--GM gives \(H_{b,\allowbreak 2}\le\allowbreak (c_*+N/2)^2\), where
\[
N=\frac m{\sqrt B}R_2,\quad
c_*=m\sqrt B-\frac{ma_1}{2\sqrt B}
\le m(1-B_*/2)\le1-\frac{1777}{3072}d\le1-37d/64.
\tag{Q58O.6.1}\label{eq:Q58O:6.1}
\]

This uses \(m\le\allowbreak 1-\kappa d/2\), \(B_*\ge\allowbreak 5d/8\), and \(d\le\allowbreak 1/3\).

\emph{The row bound.}
Write \(R_2=\allowbreak Kw\), \(w=\allowbreak g+kKR\), so
\(\|w\|_2^2=\allowbreak d-2k\rho+k^2U\).
For all \(b\ge\allowbreak 3\),
\(2m\rho-U\ge\allowbreak 7\delta_2/12\):
when \(b\ge\allowbreak 4\), use \(V\le\allowbreak \delta_2\) and \(\rho\ge\allowbreak \kappa\delta_2\);
when \(b=\allowbreak 3\), use \(\rho=\allowbreak \delta_2\) and
\(U\le\allowbreak \theta(1+\theta)\delta_2\), giving a stronger bound.
Also \(t_0-\beta_b\le\allowbreak \delta_3\le\allowbreak \theta\delta_2\) and \(t_0\le\allowbreak  d\), hence
\[
1-m^2-\gamma\ge-\beta_b^2/4-\theta \delta_2-(3/5)\beta_b d.
\]

Therefore
\[
Bd-m^2\|w\|^2\ge[7/12-(17/20+\theta)d]\delta_2
\ge79\delta_2/480\ge0.
\tag{Q58O.6.2}\label{eq:Q58O:6.2}
\]

The row Cauchy--Schwarz inequality gives \(|N|^2\le\allowbreak  d\).

\emph{The second-order filter and its coupled ratio.}
At \(z=\allowbreak i\sqrt x\), write
\(G_{b-1}(-z)=\allowbreak a_x-i\sqrt x\,\allowbreak b_x\), with
\(\kappa\le\allowbreak 1-x\le\allowbreak  a_x,\allowbreak b_x\le\allowbreak 1\).
Since \(P_1=\allowbreak 1+kz^2G_{b-1}(-z)\),
\[
|P_1|^2\le[1-kx(1-x)]^2+k^2x^3
\le(1-x+x^2)^2+x^3.
\]

The middle expression decreases for \(1\le\allowbreak  k\le\allowbreak 6/5\), because
\((6/5)x\le\allowbreak 1-x\).
Multiplying the last expression by \(x\) gives
\(x-2x^2+3x^3-x^4+x^5\), whose derivative is
\((1-2x)^2+x^2(5-4x+5x^2)>\allowbreak 0\).
Consequently
\[
ER_2^2\le(4/15)d,
\tag{Q58O.6.3}\label{eq:Q58O:6.3}
\]

The endpoint value is \(8762429/33554432<\allowbreak 4/15\), with cross-product difference \(2781293>\allowbreak 0\).

To bound the ratio \(m^2/B\), set
\(C=\allowbreak m^2B=\allowbreak m^2-mt_0+V/2\).
For \(b=\allowbreak 3,\allowbreak 4\), \(\beta_b=\allowbreak d\), \(m=\allowbreak 1-d/2\), and
\(C\ge\allowbreak  m^2-md+d^2/2\).
For \(b=\allowbreak 3\), use \(V\ge\allowbreak \delta_2\ge\allowbreak  d^2\).
For \(b=\allowbreak 4\), \(t_0=\allowbreak d-\delta_2\), so the difference between
\(-mt_0\) and \(-md\) is \(m\delta_2\ge\allowbreak  d^2/2\).
Thus
\[
C-(8/9)m^4\ge1/9-(2/9)d-d^2/12+(4/9)d^3-d^4/18
\ge79/2916>0.
\]

For \(b\ge\allowbreak 5\), \(\beta_b\le\allowbreak  d-\kappa\delta_2\le\allowbreak 77/288\), so
\(m\ge\allowbreak 499/576>\allowbreak 6/7\).
The function \(f(m)=\allowbreak 3m^2-2m-(8/9)m^4\) increases on \([6/7,\allowbreak 1]\), with
\(f(6/7)=\allowbreak 24/2401\).
If \(t_0\le\allowbreak \beta_b\), then
\(C-(8/9)m^4\ge\allowbreak  f(m)>\allowbreak 0\).
If \(t_0>\allowbreak \beta_b\), then \(b\equiv2\pmod4\), \(b\ge\allowbreak 6\), and
\(t_0-\beta_b\le\allowbreak \theta\delta_2\).
Now \(b-1\) is odd, and
\[
|G_{b-1}(i\sqrt x)|^2=\frac{1+x^{b-1}}{1+x}\ge32/45,
\]
so \(V\ge\allowbreak (32/45)\delta_2\). Therefore
\[
C-(8/9)m^4\ge24/2401-(73/1440)\delta_2
\ge24/2401-949/138240>0.
\tag{Q58O.6.4}\label{eq:Q58O:6.4}
\]

This proves \(m^2/B\le\allowbreak 9/8\) for every \(b\).
Together with \eqref{eq:Q58O:6.3}, it gives
\(\E N^2\le\allowbreak (9/8)(4/15)d=\allowbreak 3d/10\).
The signed mean is
\(\E N=\allowbreak -(md-\rho)/\sqrt B\le\allowbreak -41d/96\).
Scalar III applies to this same \(N\); using \eqref{eq:Q58O:6.1}, we obtain
\(\E H_{b,\allowbreak 2}^8\le\allowbreak 1-d/3000\).

\paragraph{Q58O.7 The ordinary-eighth catalogue}
The estimates for \(b\ge\allowbreak 3,\allowbreak L\ge\allowbreak 2\) have common deficit \(d/12000\).
Combine them with the \(b=\allowbreak 2\) positive-connector bounds of Section~\ref{proof:Q5-C050} and the unit-connector bounds of Section~\ref{proof:Q5-C051} to obtain \eqref{eq:Q58O:0.2}.
The eight-copy root construction has edge cover one, nonroot cost one and root cost \(1/8\), with the joint-law relative-entropy interface of the corresponding ordinary component.
Restore genuine equal pairs by positivity, using their fixed \(f_\ell h_\ell\le\allowbreak 1\) directions.
Complete reversal replaces \(K\) by \(-K\) and reverses the receiving arm, terminal arms, connector and equal pairs together.
